%%%%%%%%%%%%%%%%%%%%%%%%%%%%%%%%%%%%%%%%%%%%%%%%%%%%%%%%%%%%%%%%%%%%%%%%%%%%%%%%
% Preámbulo                                                                    %
%%%%%%%%%%%%%%%%%%%%%%%%%%%%%%%%%%%%%%%%%%%%%%%%%%%%%%%%%%%%%%%%%%%%%%%%%%%%%%%%

\documentclass[11pt,a4paper,titlepage,twoside,openright,openbib]{report}

%%% RELACIÓN DE VARIABLES A PERSONALIZAR %%%
\def\lingua{esp} % descomenta esta liña se redactarás a memoria en español
\def\nome{Alonso Rodriguez Iglesias}                             % substitúe aquí o teu nome
\def\nomedirectorA{Juan Touriño Domínguez}             % substitúe aquí o nome de quen dirixe
\def\nomedirectorB{María José Martín Santamaría}             % substitúe aquí o nome de quen dirixe
\def\titulo{Diseño e implementación hardware y software de un mini-supercomputador con Raspberry Pi} % substitúe aquí o título do teu TFG
\def\mencion{ENXEÑARÍA DE COMPUTADORES}

%\def\renomearcadros{si} % descomenta esta liña se redactas a memoria en español e prefires que
                         % os "cuadros" e o "índice de cuadros" se renomeen
                         % a "tablas" e "índice de tablas" respectivamente

\usepackage{estilo_tfg}

% Lista de paquetes potencialmente interesantes (uso baixo demanda)

\usepackage[all]{nowidow}
% \usepackage{alltt}       % proporciona o entorno alltt, semellante a verbatim pero que respecta comandos
% \usepackage{enumitem}    % permite personalizar os entornos de lista
% \usepackage{eurofont}    % proporciona o comando \euro
% \usepackage{float}       % permite máis opcións para controlar obxectos flotantes (táboas, figuras)
% \usepackage{hhline}      % permie personalizar as liñas horizontais en arrays e táboas
% \usepackage{longtable}   % permite construir táboas que ocupan máis dunha páxina
% \usepackage{lscape}      % permite colocar partes do documento en orientación apaisada
% \usepackage{moreverb}    % permite personalizar o entorno verbatim
% \usepackage{multirow}    % permite crear celdas que ocupan varias filas da mesma táboa
% \usepackage{pdfpages}    % permite insertar ficheiros en PDF no documento
% \usepackage{rotating}    % permite diferentes tipos de rotacións para figuras e táboas
% \usepackage{subcaption}  % permite a inclusión de varias subfiguras nunha figura
% \usepackage{tabu}        % permite táboas flexibles
% \usepackage{tabularx}    % permite táboas con columnas de anchura determinada

\usepackage{sourcecodepro}

%%%%%%%%%%%%%%%%%%%%%%%%%%%%%%%%%%%%%%%%%%%%%%%%%%%%%%%%%%%%%%%%%%%%%%%%%%%%%%%%
% Corpo                                                                        %
%%%%%%%%%%%%%%%%%%%%%%%%%%%%%%%%%%%%%%%%%%%%%%%%%%%%%%%%%%%%%%%%%%%%%%%%%%%%%%%%

\begin{document}

%%%%%%%%%%%%%%%%%%%%%%%%%%%%%%%%%%%%%%%%
% Preliminares do documento            %
%%%%%%%%%%%%%%%%%%%%%%%%%%%%%%%%%%%%%%%%

\begin{titlepage}
  
  \hspace*{128pt}
  \textcolor{udcpink}{{\fontencoding{T1}\fontfamily{phv}\selectfont Facultade de Informática}}\\[-32pt]

  \begin{center}
    \includegraphics[scale=0.3]{imaxes/udc}\\[35pt]

    {\large TRABALLO FIN DE GRAO \\
            GRAO EN ENXEÑARÍA INFORMÁTICA \\
            MENCIÓN EN \mencion } \\[100pt]
    
    \begin{huge}
      \begin{spacing}{1.3}
        \bfseries \titulo
      \end{spacing}
    \end{huge}
  \end{center}
  
  \vfill
  
  \begin{flushright}
    {\large
    \begin{tabular}{ll}
      {\bf Estudante:} & \nome \\
      {\bf Dirección:} & \nomedirectorA \\ % COPIA E PEGA ESTA LIÑA MÁIS VECES SE O PRECISAS
      {\bf Dirección:} & \nomedirectorB \\ % COPIA E PEGA ESTA LIÑA MÁIS VECES SE O PRECISAS
    \end{tabular}}
  \end{flushright}
  \rightline{A Coruña, \today.}
\end{titlepage}

\paxinaenbranco
\dedicatoria{Dedicatoria} % escribe neste comando o teu texto de dedicatoria
\paxinaenbranco
\paxinaenbranco
\begin{agradecementos}
\blindtext                % substitúe este comando polo teu texto de agradecementos
\end{agradecementos}
\paxinaenbranco
%%%%%%%%%%%%%%%%%%%%%%%%%%%%%%%%%%%%%%%%%%%%%%%%%%%%%%%%%%%%%%%%%%%%%%%%%%%%%%%%

\begin{abstract}\thispagestyle{empty}
  \blindtext % substitúe este comando polo resumo do teu TFG
             % na lingua principal do documento (tipicamente: galego)

  \vspace*{25pt}
  \begin{segundoresumo}
    \blindtext % substitúe este comando polo resumo do teu TFG
               % na lingua secundaria do documento (tipicamente: inglés)
  \end{segundoresumo}
\vspace*{25pt}
\begin{multicols}{2}
\begin{description}
\item [\palabraschaveprincipal:] \mbox{} \\[-20pt]
  \blindlist{itemize}[7] % substitúe este comando por un itemize
                         % que relacione as palabras chave
                         % que mellor identifiquen o teu TFG
                         % no idioma principal da memoria (tipicamente: galego)
\end{description}
\begin{description}
\item [\palabraschavesecundaria:] \mbox{} \\[-20pt]
  \blindlist{itemize}[7] % substitúe este comando por un itemize
                         % que relacione as palabras chave
                         % que mellor identifiquen o teu TFG
                         % no idioma secundario da memoria (tipicamente: inglés)
\end{description}
\end{multicols}

\end{abstract}

%%%%%%%%%%%%%%%%%%%%%%%%%%%%%%%%%%%%%%%%%%%%%%%%%%%%%%%%%%%%%%%%%%%%%%%%%%%%%%%%

\paxinaenbranco

\pagenumbering{roman}
\setcounter{page}{1}
\bstctlcite{IEEEexample:BSTcontrol}

\tableofcontents
\listoffigures
\listoftables
\cleardoublepage

\pagenumbering{arabic}
\setcounter{page}{1}

%%%%%%%%%%%%%%%%%%%%%%%%%%%%%%%%%%%%%%%%
% Capítulos                            %
%%%%%%%%%%%%%%%%%%%%%%%%%%%%%%%%%%%%%%%%

\chapter{Introdución}
\label{chap:introducion}

\lettrine{P}{rimer} capítulo da memoria, onde xeralmente se exporán as
liñas mestras do traballo, os obxectivos, etc. Incluimos un par de
exemplos de citas~\cite{ErlangBook,ErlangWebBook} e de referencias
internas (sección \ref{sec:mostra}, páxina \pageref{sec:mostra}).

\Blindtext

\section{Sección de mostra}
\label{sec:mostra}

\Blindtext

\subsection{Subsección de mostra}

\Blindtext

%\include{contido/...}
% Analisis -> Diseño -> Implementación -> Testeo
\include{contido/analisis_diseño}
\chapter{Implementación}
\label{chap:implementacion}

\lettrine{I}{mplementación} test.

\section{Infraestructura Hardware}
\label{sec:impl_infra_hardware}

\section{Infraestructura Software}
\label{sec:impl_infra_software}
\chapter{Medida del rendimiento}
\label{chap:medida_rendimiento}

\lettrine{L}{a} medida del rendimiento y el resultado de la misma es una parte fundamental en la evaluación y relevancia de un supercomputador, hasta el punto en el que se compite por ver cuales desarrollan un mejor resultado\footnote{\url{https://www.top500.org}}.

En este caso, si bien el rendimiento no es una prioridad, es conveniente realizar estas pruebas, especialmente para poder observar el impacto que tiene la red de comunicaciones entre los núcleos de una sola CPU, o entre múltiples CPUs y memorias. En resumen, medir la escalabilidad, que como se observa en \nameref{sec:comparacion_resultados} no es en absoluto espectacular.

\section{NAS Parallel Benchmarks}
Los \acrlong{npb} (\acrshort{npb})\footnote{\url{https://www.nas.nasa.gov/software/npb.html}} son una \textit{suite} de tests diseñados por la División de Supercomputación de la NASA para la medida del rendimiento de supercomputadores paralelos.

Estos benchmarks se dividen principalmente ocho benchmarks, teniendo otras adaptaciones para diferentes paradigmas de computación paralela, como puede ser de paralelización híbrida, computación desestructurada, o mallas de computación.

A su vez cada benchmark tiene múltiples clases, siendo estas \texttt{S, W, A, B, C, D, E}, y representando cada una un nivel de requerimientos superior, respectivamente. Específicamente, la las clases A, B y C suponen cuadruplicar el tamaño con respecto a la anterior, y las D, E y F (para los benchmarks que tienen estas clases disponibles), multiplican por 16 el tamaño del anterior.

En este caso se empleará la clase C para todos los benchmarks, excepto en MG que será clase B, y FT, que será clase A debido a limitaciones de memoria en los nodos.

\subsection{Clases y procesos}
Los benchmarks a ejecutar serán \texttt{LU, CG, FT, IS, MG y EP}, es decir, se excluyen \texttt{BT, SP y DT}. Esto es así debido a que los dos primeros requieren un número cuadrado perfecto de procesos, por lo que solo se podría ejecutar en 1, 4, 9, 16\dots procesos, y medir su rendimiento sería extraño debido a lo desbalanceado del mismo. En cuanto a DT, requiere una cantidad mínima de procesos en función de la clase del problema, procesos que para clase A y superior no son suficientes, por lo que se ha optado por descartarlo.

En concreto, los requerimientos para el número de procesos son los siguientes\footnote{Contenidos del fichero \texttt{README.install} en \texttt{NPB3.4-MPI}}:

\begin{lstlisting}
Different benchmark has different requirement for process count,
as listed below:

  BT, SP         - a square number of processes (1, 4, 9, ...)
  LU             - 2D (n1 * n2) process grid where n1/2 <= n2 <= n1
  CG, FT, IS, MG - a power-of-two number of processes (1, 2, 4, ...)
  EP, DT         - no special requirement
\end{lstlisting}

\subsection{Metodología}
La metodología es una muy importante parte en cualquier medida y análisis de la realidad, ya que permite filtrar y categorizar datos como más o menos fidedignos. En este caso la metodología a seguir ha sido moderadamente sencilla.

Se han realizado 5 ejecuciones de todos los problemas clase C, y 9 ejecuciones de los clase A o B. Esto es así, ya que a menor tiempo de ejecución, mayor probabilidad de obtener un dato atípico. Por suerte, y muy probablemente debido a la instalación mínima del sistema operativo, no ha habido variaciones significativas entre ejecuciones.

Los resultados de los benchmarks pueden encontrarse anexados en TODO \#REF\# VALUES.

\subsection{Prerrequisitos}
Para poder realizar las medidas de rendimiento, primero se necesitan tanto las herramientas de compilación como las de ejecución. Estas dependencias se deben satisfacer ejecutando:

\begin{lstlisting}[language=bash]
# Se añade la opción --needed debido a que hay paquetes del metapaquete base-devel que ya están instalados, y no es necesario reinstalar.
# Tras ejecutar el comando se aceptan las opciones por defecto
pacman -S base-devel openmpi gcc-fortran wget --needed
\end{lstlisting}

\subsection{Instalación de los \acrshort{npb}}
A continuación se deben descargar, configurar, y compilar los \acrlong{npb} únicamente en el nodo maestro. Esto se realizará como \texttt{mpiuser} en la carpeta \texttt{/mpishared} de la siguiente forma:

\begin{lstlisting}[language=bash]
# Descarga de la última versión de NPB (3.4.2)
wget https://www.nas.nasa.gov/assets/npb/NPB3.4.2.tar.gz

# Descomprimir el archivo descargado
tar xzfv NPB3.4.2.tar.gz

# Y entrar en la versión de los NPB con MPI
cd /mpishared/NPB3.4.2/NPB3.4-MPI

# Copiar los archivos de configuración para la compilación
#  1- Para el archivo make.def se añaden un par de argumentos a los flags de compilación:
#   1.1- -fallow-argument-mismatch Esto permite que el código compile en las últimas versiones de gfortran
#   1.2- -march=native Optimiza para la arquitectura específica de rpi4
sed 's/^FFLAGS\s*=\s*-O3/FFLAGS  = -O3 -fallow-argument-mismatch -march=native/g' config/make.def.template | sed -e 's/^CFLAGS\s*=\s*-O3/CFLAGS  = -O3 -march=native/g' > config/make.def

#  2- Para el archivo suite.def se hace que todos los tests sean clase C menos FT, que por razones de memoria es clase A, y MG clase B, por los mismos motivos.
sed 's/S$/C/g' config/suite.def.template | sed -e 's/ft\s*C$/ft\tA/g' | sed -e 's/mg\s*C$/mg\tB/g' > config/suite.def

# Finalmente se compilan todos los tests con make suite (con todas las CPU en paralelo)

# TODO REMOVE CLOSING DOLLAR FOR LATEX SYNTAX $
make suite -j4
\end{lstlisting}

\subsection{Preparativos para la ejecución de los benchmarks}
Para la ejecución de los benchmarks de forma paralela se necesita algún método para coordinar la ejecución del programa en todos los nodos.

\begin{itemize}
    \item Para un solo nodo con múltiples núcleos es sencillo, ya que comparten memoria, usuarios y almacenamiento
    \item Para múltiples nodos sin embargo es algo más complejo. Como ya se menciona en las Subsecciones \ref{ssec:gen_ssh} y \ref{ssec:serv_nfs}, es necesario disponer de:
    \begin{itemize}
        \item Almacenamiento compartido en forma de NFS para alojar los ejecutables
        \item Autenticación y gestión común para el lanzamiento de los programas paralelos en forma de SSH con par de claves
    \end{itemize}
\end{itemize}

A MPI se le indica qué nodos deben ejecutar las tareas mediante un archivo de hosts, o \textit{hostsfile}\cite{mpi_hostfile_option}. Este fichero contiene las direcciones de los diferentes computadores que componen el cluster, así como el número de tareas que pueden ejecutar (habitualmente coincidente con su número de núcleos)

Así, crearemos en \texttt{/mpishared} el fichero \texttt{hostfile} con los hostname de todos los nodos:

\begin{lstlisting}[language=bash]
rpi0 slots=4
rpi1 slots=4
rpi2 slots=4
rpi3 slots=4
rpi4 slots=4
rpi5 slots=4
rpi6 slots=4
rpi7 slots=4
\end{lstlisting}

\subsection{Ejecución de los benchmarks}
Ejecutar los benchmarks es moderadamente sencillo: \texttt{mpirun} funciona de la siguiente forma <TODO (posiblemente sección MPI) insertar explicación detallada> (posiblemente reubicar)

La ejecución de cada uno de los benchmarks se realizará mediante el siguiente comando \textbf{y como usuario mpiuser}, en el que deberemos especificar el número de procesos que queremos crear.

\begin{lstlisting}[language=bash]
mpirun -np <NUM_PROCS> --hostfile /mpishared/hostfile --mca opal_warn_on_missing_libcuda 0 /mpishared/NPB3.4.2/NPB3.4-MPI/bin/<KERNEL>.<CLASS>.x
\end{lstlisting}

Como estos parámetros son fijos, podemos scriptear la ejecución de los benchmarks, así como la recolección de resultados:

\begin{lstlisting}[language=bash]
#!/bin/sh

mkdir -p "results"

BENCHMARKS=( lu cg ft is mg ep )
MPIRUNFLAGS="--mca opal_warn_on_missing_libcuda 0"

for i in "${BENCHMARKS[@]}"
do
    for j in 1 2 4 8 16 32
    do
        mpirun -np $j --hostfile /mpishared/hostfile $MPIRUNFLAGS /mpishared/NPB3.4.2/NPB3.4-MPI/bin/${i}.*.x | tee -a results/$i.$j.run
    done
done
\end{lstlisting}

\section{Comparación de Resultados}
\label{sec:comparacion_resultados}
TODO
\chapter{Contenido didáctico}
\label{chap:contenido_didactico}

\lettrine{U}{na} importante parte de este proyecto es la forma en la que sirve para explicar el funcionamiento del mismo, así como realizar demostraciones prácticas de dicho funcionamiento.

Para este cometido se desarrolla una aplicación web desde la cual se podrán ejecutar ....

\section{Aplicación de monitorización}
\subsection{}
\chapter{Contido demostrativo}
\label{chap:demo}

\lettrine{E}{ntre} a introdución e as conclusións, o documento conterá
tantos capítulos como sexa preciso, sempre con coidado de non rebasar
o límite de 80 páxinas fixado polo regulamento de TFGs.

Empregaremos éste de xeito demostrativo, para ilustrar o uso de
elementos habituais que poidan ser de utilidade\footnote{Por exemplo,
  isto é unha nota a pé de páxina.}.

\section{Inclusión de imaxes}

Se precisamos imaxes no noso documento, incluirémolas do xeito que se
indica na figura~\ref{fig:exemplo} (páxina~\pageref{fig:exemplo}). Se
o facemos así, \LaTeX ubicará cada imaxe no mellor lugar posible,
lugar que pode variar a medida que o documento vaia crecendo coa
inclusión de máis texto e outros elementos (máis imaxes, táboas,
etc.).

\begin{figure}[hp!]
  \centering
  \includegraphics[width=0.75\textwidth]{imaxes/udc.png}
  \caption{Pé de imaxe descritivo}
  \label{fig:exemplo}
\end{figure}

Recoméndase almacenar os ficheiros gráficos no directorio
\texttt{imaxes}.

\subsection{Inclusión de varias sub-imaxes}

Se precisamos inserir imaxes relacionadas, pode ser apropiado
incluílas como sub-figuras, do xeito que se pode apreciar na
figura~\ref{fig:exemplo-subfiguras} coas
imaxes~\ref{fig:subfigura-rotada}
e~\ref{fig:subfigura-deformada}. Como se pode ver nos exemplos desta
sección, sempre é recomendable referirse ás imaxes pola súa
referencia, xa que dese xeito non dependemos de onde queden ubicados
os elementos en cuestión.

\begin{figure}[hp!]
  \centering
  \begin{subfigure}[c]{0.3\textwidth}
    \includegraphics[angle=45,width=\textwidth]{imaxes/udc.png}
    \caption{Pé de subimaxe rotada}
    \label{fig:subfigura-rotada}
  \end{subfigure}
  \hspace{0.1\textwidth}
  \begin{subfigure}[c]{0.3\textwidth}
    \includegraphics[width=\textwidth,height=3cm]{imaxes/udc.png}
    \caption{Pé de subimaxe deformada}
    \label{fig:subfigura-deformada}
  \end{subfigure}
  \caption{Pé de imaxe xeral}
  \label{fig:exemplo-subfiguras}
\end{figure}

\section{Inclusión de código fonte}

Se precisamos incluír fragmentos de código fonte, podemos facelo da
seguinte maneira:

\begin{lstlisting}[language=C]
#include <stdio.h>
#define N 10

int main()
{
  int i;

  // Isto é un comentario
  puts("Ola, mundo!");

  for (i = 0; i < N; i++)
  {
    puts("LaTeX é a ferramenta de edición ideal para profesionais da informática!");
  }

  return 0;
}
\end{lstlisting}

\section{Uso da relación de acrónimos e do glosario}

Os acrónimos edítanse no ficheiro \texttt{bibliografia/acronimos.tex}
e úsanse empregando a orde \texttt{acrlong} para obter o termo
completo (deste xeito: \acrlong{erlang}), a orde \texttt{acrshort}
para obter o acrónimo (deste xeito: \acrshort{erlang}). A primeira vez
que usamos un termo con acrónimo no documento é recomendable usar orde
\texttt{acrfull} (que produce ambas versións á vez:
\acrfull{erlang}). Os acrónimos que non se usan no documento, non
aparecen na relación que se xerar na versión PDF.

Pola súa banda, os termos do glosario edítanse no ficheiro
\texttt{bibliografia/glo\-sa\-rio.tex} e úsanse empregando a orde
\texttt{gls} (deste xeito, \gls{bytecode}) ou \texttt{Gls} (deste
xeito, \Gls{bytecode}). Ao igual que os acrónimos, os termos que non
se usan no documento, non aparecen na relación que se xera na versión
PDF.

\chapter{Conclusións}
\label{chap:conclusions}

\lettrine{D}{erradeiro} capítulo da memoria, onde se presentará a
situación final do traballo, as leccións aprendidas, a relación coas
competencias da titulación en xeral e a mención en particular,
posibles liñas futuras,\dots

\Blindtext


%%%%%%%%%%%%%%%%%%%%%%%%%%%%%%%%%%%%%%%%
% Apéndices, glosarios e bibliografía  %
%%%%%%%%%%%%%%%%%%%%%%%%%%%%%%%%%%%%%%%%

\appendix
\appendixpage
\chapter{Material adicional}
\label{chap:adicional}

\lettrine{E}{xemplo} de capítulo con formato de apéndice, onde se pode
incluír material adicional que non teña cabida no corpo principal do
documento, suxeito á limitación de 80 páxinas establecida no
regulamento de TFGs.

\Blindtext

%\include{anexos/...}

\printglossary[type=\acronymtype,title=\nomeglosarioacronimos]
\printglossary[title=\nomeglosariotermos]

\bibliographystyle{IEEEtranN}
\bibliography{\bibconfig,bibliografia/bibliografia}
\cleardoublepage

\end{document}

%%%%%%%%%%%%%%%%%%%%%%%%%%%%%%%%%%%%%%%%%%%%%%%%%%%%%%%%%%%%%%%%%%%%%%%%%%%%%%%%
